% ============================================================
%  Chapter 6 — Robots at Home
%  Robot World: Meet the Machines That Help Us
% ============================================================

\chapter{Robots at Home}

\chapteropening{%
  You might already live with a robot and not even realise it.
  From vacuum cleaners to lawn mowers, robots are quietly helping
  families every day.
}
\chapterbadge{Echo's Domain: The Smart Home}

% --- Introduction ---
\section*{Robots Under Your Roof}

Factory robots are big and expensive. Home robots are small, affordable,
and designed to be safe around children and pets.

\character{Echo}{I listen for your voice and control the smart lights,
  play music, and even remind you about homework. I am not just a
  speaker \textemdash{} I am a helper!}

\sectionrule

% --- Types of home robots ---
\section*{Types of Home Robots}

\subsection*{Robot Vacuums}

These disc-shaped \term{robot vacuums} zip around your floors, sucking up dust and crumbs.
They use \term{bump sensors}, \term{cliff sensors} (so they do not fall downstairs),
and sometimes \term{lidar} to map your rooms.

\begin{funfact}
  The first Roomba was sold in 2002. Since then, over 40 million
  robot vacuums have been sold worldwide.
\end{funfact}

\subsection*{\term{Robot Lawn Mowers}}

Like a vacuum for your garden. They follow a boundary wire buried in the
ground and mow a little bit every day, keeping the grass perfectly short
without you lifting a finger.

\subsection*{\term{Robot Window Cleaners}}

These cling to glass using suction and crawl across windows, spraying and
wiping as they go \textemdash{} perfect for hard-to-reach places.

\subsection*{\term{Smart Speakers} and Assistants}

Devices like Alexa, Google Home, and Siri are not physical robots, but
they use the same sense\,--\,think\,--\,act pattern: microphone (sense),
AI brain (think), speaker output and smart-home commands (act).

\subsection*{\term{Companion Robots}}

Some robots are built purely to be friends. They chat, play games,
and even recognise your face and mood.

\begin{picturethis}
  Imagine a pet that never needs feeding, never makes a mess, and
  always wants to play. Companion robots are a bit like that \textemdash{}
  though real pets still give the best cuddles!
\end{picturethis}

\sectionrule

% --- How a robot vacuum works ---
\section*{Inside a Robot Vacuum}

\begin{quickmap}
  \textbf{How a robot vacuum cleans a room:}
  \begin{enumerate}
    \item Leaves its \term{charging dock}.
    \item Spins its lidar sensor to build a map of the room.
    \item Plans an efficient zig-zag path.
    \item Drives forward, brushing and sucking dirt.
    \item If it senses a chair leg or wall, it adjusts course.
    \item When battery runs low, returns to the dock to recharge.
    \item Resumes cleaning where it left off.
  \end{enumerate}
\end{quickmap}

\sectionrule

% --- Diagram ---
\begin{diagram}[A robot vacuum's cleaning path: dock, map, clean, return.]
  \begin{tikzpicture}[>=Stealth]
    % Room outline
    \draw[MetalGray,line width=1.5pt,rounded corners=2pt] (0,0) rectangle (8,5);
    \node[font=\scriptsize\itshape] at (4,5.3) {Room};
    % Furniture obstacles
    \draw[fill=MetalGray!20,draw=MetalGray] (1.5,3.5) rectangle (2.5,4.5);
    \node[font=\tiny] at (2,4) {Table};
    \draw[fill=MetalGray!20,draw=MetalGray] (5.5,1) rectangle (7,2);
    \node[font=\tiny] at (6.25,1.5) {Sofa};
    % Charging dock
    \draw[fill=SunYellow!40,draw=SunYellow!70!black,rounded corners=2pt] (0.2,0.2) rectangle (1.2,0.8);
    \node[font=\tiny\bfseries] at (0.7,0.5) {Dock};
    % Zig-zag path
    \draw[RoboGreen,line width=1.2pt,->]
      (0.7,0.8) -- (0.7,2.5) -- (3.5,2.5) -- (3.5,0.5) -- (5,0.5)
      -- (5,3) -- (7.5,3) -- (7.5,4.5) -- (3,4.5) -- (3,3.2);
    % Robot position
    \fill[RoboGreen!80!black] (3,3.2) circle (0.2);
    \node[font=\tiny\bfseries,right] at (3.3,3.2) {Robot};
  \end{tikzpicture}
\end{diagram}

\sectionrule

% --- Experiment ---
\begin{experiment}
  \textbf{Map your room like a robot vacuum.}
  \begin{enumerate}
    \item Blindfold yourself (or close your eyes).
    \item Starting from the door, walk slowly with arms outstretched.
    \item Every time you touch something, say what it is and mark it on paper.
    \item After five minutes, open your eyes and compare your mental map
      to the real room.
  \end{enumerate}
  \textit{A robot vacuum builds its map the same way \textemdash{} by sensing
  obstacles and remembering where they are.}
\end{experiment}

\sectionrule

% --- Safety ---
\section*{Staying Safe with Home Robots}

\begin{itemize}
  \item Keep small toys off the floor so vacuum robots do not choke on them.
  \item Never ride on a robot vacuum \textemdash{} it is not a hoverboard!
  \item Be careful with voice assistants: do not share personal information
    like passwords or your home address out loud.
\end{itemize}

\begin{bigidea}
  Home robots save time on boring chores so you have more time
  for fun, family, and creativity. But they still need humans
  to set them up, maintain them, and make sure they behave.
\end{bigidea}

\sectionrule


\begin{quickcheck}
  \begin{enumerate}
    \item How does a robot vacuum know where it has already cleaned?
    \item Why must a robot vacuum return to its charging dock?
    \item Name one safety rule for living with home robots.
  \end{enumerate}
\end{quickcheck}

\sectionrule
% --- New Words ---
\begin{newwords}
  \begin{description}[labelwidth=3.8cm, leftmargin=4.0cm]
    \item[Robot vacuum] An autonomous disc-shaped robot that cleans floors.
    \item[Boundary wire] A buried wire that tells a mowing robot where to stop.
    \item[Charging dock] A station where a robot returns to refill its battery.
    \item[Companion robot] A robot designed for social interaction and company.
    \item[Smart speaker] A voice-controlled assistant with microphone and speaker.
    \item[Autonomy] The ability to act without human control.
  \end{description}
\end{newwords}
